% ============================================================================
% PLANTILLA 03 — EXAMEN DE RESOLUCIÓN DE PROBLEMAS (ciencias / matemáticas)
% ----------------------------------------------------------------------------
% Para: cálculos, aplicación de fórmulas, demostraciones matemáticas y
% problemas cuantitativos con composición matemática exigente.
% Compilar: latexmk -pdf plantilla-03-examen-problemas.tex
% Solucionario: \documentclass[soluciones]{evaluacion}
% ============================================================================
\documentclass{evaluacion}

\universidad{Universidad Nacional de San Cristóbal de Huamanga}
\facultad{Facultad de Ciencias Económicas, Administrativas y Contables}
\escuela{Escuela Profesional de Economía}
\curso{Matemática para Economistas II}
\codigocurso{EC-243}
\docente{Mg. Nombre Apellido}
\periodo{Semestre 2026-I}
\tipoevaluacion{Examen Final}
\fechaevaluacion{21 de julio de 2026}
\duracion{120 minutos}
\modalidad{Presencial · individual}

\begin{document}
\membrete

\begin{instrucciones}
\begin{itemize}
  \item Desarrolle cada problema de forma ordenada: el procedimiento vale
        tanto como el resultado. Respuestas sin justificación no reciben
        puntaje.
  \item Se permite calculadora científica no programable. El formulario
        adjunto es el único material autorizado.
  \item Exprese los resultados finales con cuatro decimales cuando
        corresponda.
\end{itemize}
\end{instrucciones}


\begin{formulario}
\begin{align*}
  \frac{d}{dx}\,a^{x} &= a^{x}\ln a
  &\int u\,dv &= uv-\int v\,du
  &\det\!\begin{pmatrix} a & b\\ c & d\end{pmatrix} &= ad-bc
\end{align*}
Condiciones de segundo orden (dos variables): $H_1=f_{xx}$,\quad
$H_2=f_{xx}f_{yy}-f_{xy}^{2}$.
\end{formulario}

% ============================================================ PROBLEMAS ====
\begin{pregunta}[5][Optimización]
Una empresa produce dos bienes en cantidades $x$ e $y$, con función de
beneficio
\[
  \pi(x,y) = 64x - 2x^{2} + 4xy - 4y^{2} + 32y - 14.
\]
\begin{apartados}
  \item Halle las condiciones de primer orden y el punto crítico. \pts{2}
  \item Verifique las condiciones de segundo orden y clasifique el punto
        crítico. \pts{2}
  \item Calcule el beneficio máximo. \pts{1}
\end{apartados}
\espaciorespuesta{8cm}
\begin{solucion}
\textbf{a)} Las condiciones de primer orden son
\[
  \pi_x = 64 - 4x + 4y = 0, \qquad
  \pi_y = 4x - 8y + 32 = 0.
\]
De la primera, $x = 16 + y$; sustituyendo en la segunda:
$4(16+y) - 8y + 32 = 0 \Rightarrow y^{*} = 24$, $x^{*} = 40$.

\textbf{b)} El hessiano es
\[
  H = \begin{pmatrix} \pi_{xx} & \pi_{xy}\\ \pi_{yx} & \pi_{yy}\end{pmatrix}
    = \begin{pmatrix} -4 & 4\\ 4 & -8\end{pmatrix},
  \qquad H_1 = -4 < 0,\quad H_2 = 32-16 = 16 > 0,
\]
por lo que $(40,24)$ es un \emph{máximo} local (y global, por concavidad
estricta de $\pi$).

\textbf{c)} $\pi(40,24) = 64(40) - 2(1600) + 4(960) - 4(576) + 32(24) - 14
= 1\,650$.
\end{solucion}
\end{pregunta}

\begin{pregunta}[5][Demostración]
Sea $f\colon\mathbb{R}^{n}_{+}\to\mathbb{R}$ una función de producción
homogénea de grado uno y diferenciable. Demuestre el teorema de Euler:
\[
  \sum_{i=1}^{n} x_i\,\frac{\partial f}{\partial x_i}(x) = f(x).
\]
Interprete económicamente el resultado bajo remuneración a productividades
marginales.
\espaciorespuesta{7cm}
\begin{solucion}
Por homogeneidad de grado uno, $f(\lambda x)=\lambda f(x)$ para todo
$\lambda>0$. Derivando ambos lados respecto de $\lambda$:
\[
  \frac{d}{d\lambda} f(\lambda x)
  = \sum_{i=1}^{n} \frac{\partial f}{\partial x_i}(\lambda x)\,x_i
  = f(x).
\]
Evaluando en $\lambda = 1$ se obtiene
$\sum_i x_i f_i(x) = f(x)$. \hfill$\blacksquare$

\emph{Interpretación:} si cada factor se paga según su productividad
marginal, el producto se agota exactamente entre los factores
(teorema de agotamiento del producto); no quedan beneficios puros bajo
rendimientos constantes a escala.
\end{solucion}
\end{pregunta}

\begin{pregunta}[5][Ecuaciones diferenciales]
El precio $p(t)$ de un mercado se ajusta según
\[
  \dot{p} = \alpha\,\bigl[D(p) - S(p)\bigr], \qquad \alpha > 0,
\]
con $D(p) = 24 - 2p$ y $S(p) = -6 + 4p$.
\begin{apartados}
  \item Halle el precio de equilibrio $p^{*}$ y resuelva la ecuación
        diferencial para $p(0)=p_0$. \pts{3}
  \item Analice la estabilidad dinámica del equilibrio. \pts{2}
\end{apartados}
\espaciorespuesta{7cm}
\begin{solucion}
\textbf{a)} En equilibrio $D=S$: $24-2p = -6+4p \Rightarrow p^{*}=5$.
La ecuación se escribe $\dot p = \alpha(30 - 6p) = -6\alpha\,(p-5)$,
ecuación lineal de primer orden cuya solución es
\[
  p(t) = 5 + (p_0 - 5)\,e^{-6\alpha t}.
\]
\textbf{b)} Como $-6\alpha<0$, $\lim_{t\to\infty}p(t)=5$ para cualquier
$p_0$: el equilibrio es \emph{globalmente estable} y la convergencia es
monótona a tasa $6\alpha$.
\end{solucion}
\end{pregunta}

\begin{pregunta}[5][Álgebra lineal]
Considere el modelo insumo-producto de Leontief con matriz de coeficientes
técnicos y vector de demanda final
\[
  A = \begin{pmatrix} 0{,}2 & 0{,}3\\ 0{,}4 & 0{,}1 \end{pmatrix},
  \qquad
  d = \begin{pmatrix} 120\\ 90 \end{pmatrix}.
\]
\begin{apartados}
  \item Verifique que la economía es productiva (condición de
        Hawkins--Simon). \pts{2}
  \item Calcule el vector de producción bruta $x = (I-A)^{-1}d$. \pts{3}
\end{apartados}
\espaciorespuesta{6.5cm}
\begin{solucion}
\textbf{a)} $I-A = \begin{pmatrix} 0{,}8 & -0{,}3\\ -0{,}4 & 0{,}9
\end{pmatrix}$. Los menores principales dominantes son $0{,}8>0$ y
$\det(I-A)=0{,}72-0{,}12=0{,}60>0$: se cumple Hawkins--Simon.

\textbf{b)}
\[
  (I-A)^{-1} = \frac{1}{0{,}60}
  \begin{pmatrix} 0{,}9 & 0{,}3\\ 0{,}4 & 0{,}8 \end{pmatrix}
  = \begin{pmatrix} 1{,}5 & 0{,}5\\ 0{,}\overline{6} & 1{,}\overline{3}
  \end{pmatrix},
  \qquad
  x = \begin{pmatrix} 1{,}5(120)+0{,}5(90)\\
        0{,}\overline{6}(120)+1{,}\overline{3}(90)\end{pmatrix}
    = \begin{pmatrix} 225\\ 200 \end{pmatrix}.
\]
\end{solucion}
\end{pregunta}

\findelexamen
\end{document}
